\section{Trabajos relacionados}

Caldarelli y Ellul revisaron propuestas que utilizan blockchain para conservar
y verificar expedientes académicos. El estudio destaca la trazabilidad de los
registros, la reducción de intermediarios y la necesidad de identificar a la
institución que introduce la información en la cadena
\cite{Caldarelli2021}. El proyecto toma como referencia la importancia de una
bitácora auditable y de la integridad documental, pero no adopta blockchain:
conserva los documentos y eventos dentro de la infraestructura institucional y
se concentra en el flujo de pre-sustentaciones de la UTEQ.

Fekete y Kiss propusieron un sistema de educación superior basado en contratos
inteligentes y una arquitectura modular con zkEVM. Su propuesta busca que los
documentos y las interacciones académicas puedan verificarse a lo largo del
tiempo, con énfasis en seguridad, descentralización e interoperabilidad
\cite{Fekete2023}. De este trabajo se toma la separación de responsabilidades y
la necesidad de registrar cada interacción relevante. La diferencia es que el
sistema de este proyecto emplea servicios web y autenticación institucional,
sin billeteras ni transacciones en una red blockchain.

Satybaldy, Subedi y Nowostawski diseñaron un marco para verificar documentos en
línea mediante identidad autosoberana y credenciales verificables. Su prueba de
concepto prioriza privacidad, interoperabilidad y verificación sin depender de
un intermediario \cite{Satybaldy2022}. El proyecto toma como referencia la
protección de los documentos y el control de acceso por actor, pero su alcance
es diferente: administra solicitudes, informes, evaluaciones y actas dentro de
un procedimiento universitario concreto, en vez de crear credenciales
portables para múltiples organizaciones.

Los tres trabajos muestran mecanismos para preservar confianza y trazabilidad
en documentos digitales. Ninguno integra, dentro del proceso específico de una
pre-sustentación de la UTEQ, la entrega del informe y del reporte de similitud,
el seguimiento de la solicitud, la evaluación del tribunal y la generación del
acta. Ese flujo institucional constituye la diferencia principal del sistema
propuesto.
